\documentclass[11pt]{article}
\usepackage[a4paper,margin=0.82in]{geometry}
\usepackage{amsmath,amssymb,amsthm,mathtools}
\usepackage{enumitem,microtype}
\usepackage[dvipsnames]{xcolor}
\usepackage[colorlinks=true,linkcolor=MidnightBlue,citecolor=MidnightBlue,urlcolor=MidnightBlue]{hyperref}
\setlist{nosep,leftmargin=*}
\newtheorem{theorem}{Theorem}
\newtheorem{lemma}[theorem]{Lemma}
\newcommand{\gra}{\operatorname{gra}}
\newcommand{\domn}{\operatorname{dom}}
\newcommand{\cl}{\operatorname{cl}}
\newcommand{\co}{\operatorname{co}}
\newcommand{\inte}{\operatorname{int}}
\title{Two positive regimes for the closed-graph/property-\((Q)\) problem\\
\large Convex graphs and relative domain interior}
\author{Autonomous solve-first investigation, lane 12}
\date{13 August 2026}

\begin{document}
\maketitle

\begin{abstract}
Borwein and Yao ask whether a maximally monotone operator with norm
$\times$ weak-star closed graph can differ from its Cesari reconstruction
$\widehat A$.  We prove two positive subcases.  First, their reflexive
closed-convex-graph example is valid on every Banach space; neither
reflexivity nor monotonicity is needed.  Second, their nonempty-domain-
interior theorem extends to operators whose domain has nonempty relative
interior in its closed affine hull.  The latter result does not require graph
closedness.  We also isolate why three natural counterexample constructions
fail: the same unbounded cancellation that enlarges the Cesari hull puts the
missing pair into the norm $\times$ weak-star graph closure.  The general
empty-relative-interior, nonconvex-graph case is not claimed.
\end{abstract}

\section{Question and notation}

Let $X$ be a real Banach space and $A:X\rightrightarrows X^*$.  Following
Borwein--Yao \cite{BY}, define
\begin{equation}\label{eq:hat}
 \widehat A(x)=\bigcap_{\varepsilon>0}
 \cl^{w^*}\co A(x+\varepsilon B_X).
\end{equation}
Their Theorem 5.3 proves $\widehat A=A$ when $A$ is maximally monotone and
$\inte\domn A\ne\varnothing$.  At the end of their Example 6.2 they ask
whether $\widehat A$ and $A$ can differ when $A$ is maximally monotone and
$\gra A$ is norm $\times$ weak-star closed.

The implications relevant here are
\[
 \widehat A=A\quad\Longrightarrow\quad A\text{ has property (Q)}
 \quad\Longrightarrow\quad \gra A\text{ is norm $\times$ weak-star closed}.
\]
The question asks whether the last implication can be reversed under maximal
monotonicity.  The following two theorems enlarge the known positive region,
but do not cover every maximally monotone graph.

\section{Convex graphs: reflexivity is unnecessary}

\begin{theorem}[closed convex graph]\label{thm:convex}
Let $X$ be any Banach space.  If $\gra A$ is nonempty, convex, and closed for
the norm $\times$ weak-star topology, then $\widehat A=A$.
No monotonicity assumption is required.
\end{theorem}

\begin{proof}
The inclusion $A\subseteq\widehat A$ is immediate.  Conversely, fix
$(x,x^*)\in\gra\widehat A$, a norm neighborhood $U$ of $x$, and a weak-star
neighborhood $V$ of $x^*$.  Choose $\varepsilon>0$ with
$x+\varepsilon B_X\subseteq U$.  By \eqref{eq:hat}, there are
$x_i\in x+\varepsilon B_X$, $x_i^*\in A(x_i)$, and convex coefficients
$\lambda_i$ such that
\[
 z^*:=\sum_i\lambda_i x_i^*\in V.
\]
Convexity of the graph gives
\[
 \left(z,z^*\right):=\left(\sum_i\lambda_i x_i,
 \sum_i\lambda_i x_i^*\right)\in\gra A.
\]
Moreover $z\in x+\varepsilon B_X\subseteq U$.  Thus every product
neighborhood $U\times V$ of $(x,x^*)$ meets $\gra A$.  Closedness of the
graph yields $(x,x^*)\in\gra A$.
\end{proof}

Borwein--Yao stated the corresponding example under reflexivity and replaced
weak-star closed convex hulls by norm-closed convex hulls.  The product-
neighborhood argument shows that convexity of the graph itself is the only
ingredient needed to turn every local convex combination into a graph point.
In particular, Theorem~\ref{thm:convex} applies to every norm $\times$
weak-star closed maximal monotone linear relation.

\section{Relative interior in the closed affine hull}

\begin{theorem}[relative-interior extension]\label{thm:relative}
Let $A:X\rightrightarrows X^*$ be maximally monotone and choose
$x_0\in\domn A$.  Put
\[
 Y=\overline{\operatorname{span}}(\domn A-x_0).
\]
If $\domn A-x_0$ has nonempty interior as a subset of the Banach space $Y$,
then $\widehat A=A$.  Consequently $A$ has property (Q) and its graph is norm
$\times$ weak-star closed.
\end{theorem}

\begin{proof}
Let $R:X^*\to Y^*$ be restriction.  Maximality first gives the saturation
identity
\begin{equation}\label{eq:saturation}
 A(x)+Y^\perp=A(x)\qquad(x\in\domn A).
\end{equation}
Indeed, if $a^*\in A(x)$ and $n^*\in Y^\perp$, then for every
$(z,z^*)\in\gra A$ we have $x-z\in Y$ and hence
\[
 \langle x-z,a^*+n^*-z^*\rangle
 =\langle x-z,a^*-z^*\rangle\ge0.
\]
Thus $(x,a^*+n^*)$ is monotonically related to $\gra A$, so maximality puts
it in the graph.

Define $B:Y\rightrightarrows Y^*$ by
\[
 B(y)=R\bigl(A(x_0+y)\bigr).
\]
This operator is maximally monotone on $Y$.  To see maximality, suppose
$(y,y^*)$ is monotonically related to $\gra B$ and extend $y^*$ to some
$x^*\in X^*$ by Hahn--Banach.  The defining inequalities then say precisely
that $(x_0+y,x^*)$ is monotonically related to $\gra A$.  Maximality of $A$
gives $x^*\in A(x_0+y)$, so $y^*\in B(y)$.

By hypothesis $\inte_Y\domn B\ne\varnothing$.  Theorem 5.3 of \cite{BY},
applied in the Banach space $Y$, yields $\widehat B=B$.

It remains to lift this identity.  If $x\notin x_0+Y$, then
$\widehat A(x)=A(x)=\varnothing$, because the closed affine space $x_0+Y$
has positive distance from $x$.  Let therefore $x=x_0+y$ and
$x^*\in\widehat A(x)$.  Restriction is weak-star to weak-star continuous,
and every domain point in $x+\varepsilon B_X$ differs from $x_0$ by an
element of $y+\varepsilon B_Y$.  Hence, for every $\varepsilon>0$,
\[
 Rx^*\in\cl^{w^*}\co B(y+\varepsilon B_Y),
\]
so $Rx^*\in\widehat B(y)=B(y)$.  Choose $a^*\in A(x)$ with
$Ra^*=Rx^*$.  Then $x^*-a^*\in Y^\perp$, and \eqref{eq:saturation} gives
$x^*\in A(x)$.  Therefore $\widehat A\subseteq A$; the reverse inclusion is
automatic.
\end{proof}

Theorem~\ref{thm:relative} includes the source theorem ($Y=X$), but also
covers operators supported on a proper closed affine subspace, where the
ordinary interior in $X$ is empty.  Notice that graph closedness is a
conclusion, not a hypothesis.

\section{Why the full question remains difficult}

For a convex combination of local graph values, monotonicity produces the
covariance identity
\begin{align*}
 &\sum_i\lambda_i\langle d_i,e_i\rangle
 -\left\langle\sum_i\lambda_i d_i,
                    \sum_i\lambda_i e_i\right\rangle\\
 &\hspace{25mm}=\frac12\sum_{i,j}\lambda_i\lambda_j
 \langle d_i-d_j,e_i-e_j\rangle\ge0.
\end{align*}
Here $d_i=x_i-x$ and $e_i=x_i^*-x^*$.  The sign is insufficient to show
that $(x,x^*)$ is monotonically related to the graph, because the local dual
values may grow faster than $\|x_i-x\|^{-1}$.

Three counterexample templates were tested in detail: a supremum-type
subdifferential on $\ell_1$, a polyhedral normal cone on $\ell_1$, and a
smooth compact ellipsoid in $\ell_2$.  Each makes $\widehat A(x)$ strictly
larger than $A(x)$ by averaging two unbounded local values.  Each fails the
closed-graph hypothesis for the same reason: after finitely many weak-star
tests are fixed, unused high-coordinate directions can be chosen so that an
\emph{actual} graph value annihilates those tests while its base point tends
to $x$.  Thus the allegedly missing pair already belongs to the product
closure of the graph.

This is evidence for a positive answer, but not a proof.  Turning the finite-
test cancellation into a general selection theorem appears to require
maximality of finite-dimensional restrictions of $A$; that maximality is
unavailable without an interior or sum-theorem qualification.  Theorem
\ref{thm:relative} is exactly the regime in which the restriction argument
can be completed.  The surviving hard core is therefore the case where the
domain has empty relative interior in its closed affine hull and the graph is
genuinely nonconvex.

\section{Verification and literature scope}

Both proofs use only product-neighborhood closure, Hahn--Banach extension,
maximality, and the source's proved nonempty-interior theorem.  No sequence is
substituted for a net and no boundedness of weak-star convergent nets is
assumed.  This last point is essential: the graph-closure pathology is driven
by unbounded nets.

A bounded search on 13 August 2026 used the exact question, the title and
citations of \cite{BY}, ``property (Q)'', and norm $\times$ weak-star closure
terminology.  It found the source's survey reprise and later work on bounded
strong $\times$ weak-star closure operations, but no paper resolving the
exact converse.  The present packet is therefore classified conservatively
as a substantial partial result, not as a solution of the full question.

\begin{thebibliography}{9}
\bibitem{BY}
J. M. Borwein and L. Yao,
\emph{Structure theory for maximally monotone operators with points of
continuity}, J. Optim. Theory Appl. 157 (2013), 1--24;
arXiv:1203.1101.
\end{thebibliography}

\end{document}
